% GENERATED FILE. Edit canonical lessons, then run npm run generate:books.
% canonical-source-hash: fnv1a64:ac25a545d0a826cd
% canonical-lessons: ZH-C341-dousha, ZH-C341-yuebing, ZH-W341-bing-cake, ZH-C341-xianr, ZH-C341-pinchang, ZH-C341-shengyi

\chapter{Mooncakes}
\label{ch:zh-mooncakes}
\hlchaptermodality{341}

\begin{chapteropening}
\textbf{By the end of this chapter:} \emph{I can talk about mooncakes, bean paste, fillings and tasting, and say business is good.}

Complete the last lesson of chapter 341: I can talk about mooncakes, bean paste, fillings and tasting, and say business is good.
\end{chapteropening}

\section[dòushā]{dòushā — sweet bean paste}

\label{lesson:ZH-C341-dousha}



\begin{quote}
Before the new one: say the Mandarin for a solar term of the year, then the Mandarin for to tend the family graves.
\end{quote}

\subsection*{You'll want to know: dòushā}
\emph{dòushā} — \textquotedblleft{}sweet bean paste\textquotedblright{}.

\begin{quote}
\emph{zhège miànbāo lǐ yǒu dòushā} — There's bean paste in this bun.
\end{quote}

Its characters come later; for now it is a word for the ear.

\begin{cousinweb}
\emph{dòushā} is built from \emph{dòu} 'beans' and \emph{shā} 'sand', both inherited from Middle Chinese, \emph{duwH} and \emph{srae} (Guangyun, in Baxter's transcription), which glosses the first 'beans' and the second 'grit scattered by water'.
\end{cousinweb}

\begin{grammarlens}[title={What you've built}]
One more word for telling what happened.
\end{grammarlens}

\subsection*{Your turn}
\begin{itemize}
\raggedright
  \item \emph{Say it:} \emph{dòushā}
  \item \emph{Say it:} \emph{dòushā}, once more
  \item \emph{Say it:} \emph{zhège miànbāo lǐ yǒu dòushā} — There's bean paste in this bun.
  \item \emph{Recall:} say the Mandarin for a solar term of the year, then the Mandarin for to tend the family graves, then say \emph{dòushā} again
\end{itemize}

\begin{tcolorbox}[breakable,colback=teal!4,colframe=teal!35!black,title={Before you move on}]
What does dòushā mean? (\textquotedblleft{}Sweet bean paste\textquotedblright{}.) Say it once more.
\end{tcolorbox}
\section[\texorpdfstring{\zh{月饼} (yuèbǐng)}{月饼 (yuèbǐng)}]{\zh{月饼} (yuèbǐng) — mooncake}

\label{lesson:ZH-C341-yuebing}



\begin{quote}
Before the new one: say the Mandarin for to tend the family graves, then the Mandarin for sweet bean paste.
\end{quote}

\subsection*{You'll want to know: \zh{月饼}}
\textbf{\zh{月饼}} — \emph{yuèbǐng} — \textquotedblleft{}mooncake\textquotedblright{}.

\begin{quote}
\textbf{\zh{我买了两个月饼}} \emph{wǒ mǎi le liǎng ge yuèbǐng} — I bought two mooncakes.
\end{quote}

The new character, in pieces: \zh{饼} — nine strokes: on the left the food piece of \zh{饭}, three strokes; on the right two short strokes, two bars, and two strokes down. The next lesson puts it in your hand.

\begin{cousinweb}
\textbf{\zh{月饼}} (\emph{yuèbǐng}) is built from \zh{月} 'the moon' and \zh{饼} 'a cake', both inherited from Middle Chinese, \emph{ngjwot} and \emph{pjiengX} (Guangyun, in Baxter's transcription), which gives the second only its reading.
\end{cousinweb}

\begin{grammarlens}[title={What you've built}]
One more word for telling what happened.
\end{grammarlens}

\subsection*{Your turn}
\begin{itemize}
\raggedright
  \item \emph{Say it:} \emph{yuèbǐng}
  \item \emph{Say it:} \emph{yuèbǐng}, once more
  \item \emph{Say it:} \emph{wǒ mǎi le liǎng ge yuèbǐng} — I bought two mooncakes.
  \item \emph{Recall:} say the Mandarin for to tend the family graves, then the Mandarin for sweet bean paste, then say \emph{yuèbǐng} again
\end{itemize}

\begin{tcolorbox}[breakable,colback=teal!4,colframe=teal!35!black,title={Before you move on}]
What does \zh{月饼} mean? (\textquotedblleft{}Mooncake\textquotedblright{}.) Say it once more.
\end{tcolorbox}
\section[\texorpdfstring{\zh{饼} (bǐng)}{饼 (bǐng)}]{\zh{饼} — the food piece, then two bars and two legs}

\label{lesson:ZH-W341-bing-cake}



\begin{quote}
The word first, then the character inside it.

\begin{itemize}
\raggedright
  \item \emph{Recall:} say \emph{yuèbǐng} — \textbf{R1}, one lesson back, before it decays
\end{itemize}

\begin{quote}
\zh{饭} \zh{月}
\end{quote}

\emph{Read it:} each one aloud, with nothing said first — \emph{fàn}, \emph{yuè}
\end{quote}

\subsection*{Script — the food piece, then two bars and two legs}
\hlblockfigure{\detokenize{figures/ZH-W341-bing-cake-filmstrip.pdf}}{How \zh{饼} is written, stroke by stroke}

\begin{quote}
\zh{饼}
\end{quote}

\textbf{Two pieces, left before right.}

\textbf{The left piece, the food piece, three strokes,} as in \zh{饭}: a stroke falling to the left, a short bar that turns down to the left, then a vertical that flicks up to the right.

\textbf{The right piece, six strokes.} Two short strokes on top, the left one a dot, the right one falling to the left. Two bars, the lower one longer. Then a stroke down the left that falls away to the left, and a vertical down the right.

3 + 6 = nine.

\subsection*{Your turn}
\begin{itemize}
\raggedright
  \item \emph{Write it:} \zh{饼}, all nine, with the model in front of you
  \item \emph{Cover:} write it again from memory, and count the pieces, not the strokes
\end{itemize}

\begin{tcolorbox}[breakable,colback=teal!4,colframe=teal!35!black,title={Before you move on}]
How many strokes in \zh{饼}? (\textbf{Nine}.) Where have you met the left piece? (\textbf{On the left of \zh{饭}}.) How many bars on the right? (\textbf{Two}.)

Source: \href{https://raw.githubusercontent.com/chanind/hanzi-writer-data/68d10a4b21150cae5e1ebbd223eed289cf32d90c/data/\%E9\%A5\%BC.json}{Hanzi Writer Data \zh{饼}, pinned stroke paths and medians}.
\end{tcolorbox}
\section[xiànr]{xiànr — filling}

\label{lesson:ZH-C341-xianr}



\begin{quote}
Before the new one: say the Mandarin for sweet bean paste, then the Mandarin for a mooncake.
\end{quote}

\subsection*{You'll want to know: xiànr}
\emph{xiànr} — \textquotedblleft{}filling\textquotedblright{}.

\begin{quote}
\emph{zhège yuèbǐng de xiànr shì dòushā} — This mooncake has a bean-paste filling.
\end{quote}

Its characters come later; for now it is a word for the ear.

\begin{cousinweb}
\emph{xiànr} No dictionary we use records an older form, so its history is not settled. The ending r is the \zh{儿} of \zh{一会儿}.
\end{cousinweb}

\begin{grammarlens}[title={What you've built}]
One more word for telling what happened.
\end{grammarlens}

\subsection*{Your turn}
\begin{itemize}
\raggedright
  \item \emph{Say it:} \emph{xiànr}
  \item \emph{Say it:} \emph{xiànr}, once more
  \item \emph{Say it:} \emph{zhège yuèbǐng de xiànr shì dòushā} — This mooncake has a bean-paste filling.
  \item \emph{Recall:} say the Mandarin for sweet bean paste, then the Mandarin for a mooncake, then say \emph{xiànr} again
\end{itemize}

\begin{tcolorbox}[breakable,colback=teal!4,colframe=teal!35!black,title={Before you move on}]
What does xiànr mean? (\textquotedblleft{}Filling\textquotedblright{}.) Say it once more.
\end{tcolorbox}
\section[pǐncháng]{pǐncháng — to taste; to sample}

\label{lesson:ZH-C341-pinchang}



\begin{quote}
Before the new one: say the Mandarin for a mooncake, then the Mandarin for the filling of a dumpling.
\end{quote}

\subsection*{You'll want to know: pǐncháng}
\emph{pǐncháng} — \textquotedblleft{}to taste; to sample\textquotedblright{}.

\begin{quote}
\emph{dàjiā dōu pǐncháng le wǒ zuò de yuèbǐng} — Everyone tasted the mooncakes I made.
\end{quote}

Its characters come later; for now it is a word for the ear.

\begin{cousinweb}
\emph{pǐncháng} is built from \emph{pǐn} and \emph{cháng}, both inherited from Middle Chinese, \emph{phimX} and \emph{dzyang} (Guangyun, in Baxter's transcription), which glosses the first 'a rank' and 'a kind', and the second 'to try' and, quoting an older dictionary, 'to taste with the mouth'.
\end{cousinweb}

\begin{grammarlens}[title={What you've built}]
One more word for telling what happened.
\end{grammarlens}

\subsection*{Your turn}
\begin{itemize}
\raggedright
  \item \emph{Say it:} \emph{pǐncháng}
  \item \emph{Say it:} \emph{pǐncháng}, once more
  \item \emph{Say it:} \emph{dàjiā dōu pǐncháng le wǒ zuò de yuèbǐng} — Everyone tasted the mooncakes I made.
  \item \emph{Recall:} say the Mandarin for a mooncake, then the Mandarin for the filling of a dumpling, then say \emph{pǐncháng} again
\end{itemize}

\begin{tcolorbox}[breakable,colback=teal!4,colframe=teal!35!black,title={Before you move on}]
What does pǐncháng mean? (\textquotedblleft{}To taste; to sample\textquotedblright{}.) Say it once more.
\end{tcolorbox}
\section[\texorpdfstring{\zh{生意} (shēngyi)}{生意 (shēngyi)}]{\zh{生意} (shēngyi) — business; trade}

\label{lesson:ZH-C341-shengyi}



\begin{quote}
Before the new one: say the Mandarin for the filling of a dumpling, then the Mandarin for to taste food with care.
\end{quote}

\subsection*{You'll want to know: \zh{生意}}
\textbf{\zh{生意}} — \emph{shēngyi} — \textquotedblleft{}business; trade\textquotedblright{}.

\begin{quote}
\textbf{\zh{每年这个时候月饼的生意最好}} \emph{měinián zhège shíhou yuèbǐng de shēngyi zuì hǎo} — Every year at this time the mooncake trade is at its best.
\end{quote}

\begin{cousinweb}
\textbf{\zh{生意}} (\emph{shēngyi}) is built from \zh{生} 'life' and \zh{意} 'intent', both inherited from Middle Chinese, \emph{srjaeng} and \emph{'iH} (Guangyun, in Baxter's transcription), which glosses the first 'to grow' and the second 'the will'. How the pair came to mean business is not settled.
\end{cousinweb}

\begin{grammarlens}[title={What you've built}]
One more word for telling what happened.
\end{grammarlens}

\subsection*{Your turn}
\begin{itemize}
\raggedright
  \item \emph{Say it:} \emph{shēngyi}
  \item \emph{Say it:} \emph{shēngyi}, once more
  \item \emph{Say it:} \emph{měinián zhège shíhou yuèbǐng de shēngyi zuì hǎo} — Every year at this time the mooncake trade is at its best.
  \item \emph{Recall:} say the Mandarin for the filling of a dumpling, then the Mandarin for to taste food with care, then say \emph{shēngyi} again
\end{itemize}

\begin{tcolorbox}[breakable,colback=teal!4,colframe=teal!35!black,title={Before you move on}]
What does \zh{生意} mean? (\textquotedblleft{}Business; trade\textquotedblright{}.) Say it once more.
\end{tcolorbox}
